The OpenMC model uses a fixed-source problem on a three-gram spherical sample of a fissile target.
A uniformly distributed source is placed inside the target with a vacuum boundary condition at the edge of the sphere.
The sample must have a sufficiently large volume to ensure a high number of fissions, but excessive volume increases self-shielding and self-multiplication.
Adjusting the density mitigates this issue: higher density increases self-shielding and self-multiplication, while lower density reduces fission probability and worsens statistical accuracy.

The fixed-source in this model consists entirely of neutrons at a specific energy, rather than a spectrum.
This energy dictates fission cross sections, transmutation cross sections, and fission yields.
Traditionally, tabulated group parameter data is formatted for a single energy or energy range, such as fast or thermal.
However, macroscopic experiments, discussed in Section \ref{sec:measurementofdnps}, use reactors with varying neutron energy spectra, not mono-energetic neutrons.
This could be addressed by modeling a reactor in OpenMC and using a k-eigenvalue solve to generate a neutron energy spectrum.
However, this work assumes the spectral effects are sufficiently small to exclude.
Photons can also generate delayed neutrons through photofission, though this is not investigated here \cite{kinlaw2007delayed}.

The neutron source distribution does not replicate the experimental process, where neutrons irradiate the sample externally rather than internally.
However, an evenly distributed source eliminates self-shielding concerns.
The exact spatial locations of fission events within the sample should negligibly affect the nuclide concentrations, which are the outputs of the OpenMC simulation.
With a density of 10 $\frac{g}{cc}$, the spherical sample has a radius of approximately 0.4 $cm$.
The spatial resolution of concentrations within the sample is insignificant.
Additionally, the concentrations extracted from the OpenMC simulation are spatially collapsed in the post-processing, so the spatial effect of a uniformly distributed source serves only to reduce any self-shielding effects which could reduce the statistics and cause more transmutation than expected.

After running the OpenMC simulation, the concentrations are collected from the depletion results for the final irradiation time step.
These final concentrations are then converted into a delayed neutron count rate.